\section{The one-value endpoint multiplier}\label{app:optimization}

The endpoint $19m-1$ can be sharpened slightly. We first give
an explicit refinement and then compare its multiplier with
the transition of the decay estimate.

\subsection{A sharper uniform bound}

\begin{proposition}\label{prop:one-sharp}
For every positive odd integer $m$, put
\begin{equation}\label{eq:one-endpoint}
 s(m)=m-1+2\left\lceil\frac{879m-50}{100}\right\rceil<18.58m.
\end{equation}
At least one of $\beta(m+1),\beta(m+3),\ldots,\beta(s(m))$
is irrational.
\end{proposition}

\begin{proof}
For $B=1$, put
\[
 a=2\left\lceil\frac{879m-50}{100}\right\rceil.
\]
Then $a\ge(879/50)m-1$ and $a/m\ge52/3$.
Indeed, $m=1$ gives $a=18$, $m=3$ gives $a=52$, and
$m\ge5$ gives $a/m\ge879/50-1/5>52/3$.
With $r=2$ and $u=1/100$, the comparison
\begin{equation}\label{eq:f-sharp}
 1001^3\left(\frac{401}{601}\right)^{52}<1
\end{equation}
implies $f(u)<1$ for every $m$.

With $e_*$ as in \eqref{eq:e-tight}, put
$A_*=\log(27/(4e_*))$ and $C=5\log(20.02)$.
Then
\begin{align*}
 \frac12\log U
 &=mC-A_*a+(m-1)\log e_*\\
 &\le m\left(C+\log e_*-\frac{879}{50}A_*\right)
                                      +A_*-\log e_*.
\end{align*}
Both terms on the last line are negative. For the first,
take logarithms of the exact rational comparison
\begin{equation}\label{eq:sharp-U}
 e_*^{929}2^{250}
       \left(\frac{1001}{100}\right)^{250}
       \left(\frac4{27}\right)^{879}<1.
\end{equation}
For the second, use $27/4<(163/60)^2<e_*^2$.
The criterion gives the stated interval. Finally,
$\lceil x\rceil<x+1$ gives $s(m)<18.58m$.
\end{proof}

\subsection{The limiting multiplier}

The multiplier $18.58$ comes from a rational estimate close
to the transition of the decay bound. The limiting profile
also explains the choice $r=2$.

For $B=1$, write $\mu=a/m$ and $y=x/2$. Since $D=1$,
$K=2^L$ and $k=m-1$, the profile in \eqref{eq:Phi} satisfies
the exact identity
\begin{equation}\label{eq:finite-profile}
 a+k+\Phi(2y)=m\Psi_r(\mu,y)-1,
\end{equation}
where
\begin{equation}\label{eq:limiting-profile}
 \begin{split}
 \Psi_r(\mu,y)={}&\mu+1+2L\log2
       +(y+L)\log(y+L)-y\log y\\
 &-\mu\bigl((y+r+1)\log(y+r+1)
                         -(y+r)\log(y+r)\bigr),
 \qquad L=2r+1.
 \end{split}
\end{equation}
The constant $+1$ in $\Psi_r$ comes from the additional denominator cost of
$m-1$ derivatives; the final $-1$ in \eqref{eq:finite-profile}
accounts for the difference between $m-1$ and $m$.
Taking the maximum gives
$\log(\ep^{a+k}G)=m\max_{y\ge0}\Psi_r(\mu,y)-1$.
After division by $m$, the correction is $-1/m$ and vanishes
in the limit. Decay in the limiting profile corresponds to
$\max_{y\ge0}\Psi_r(\mu,y)<0$.
For $r=2$, numerical solution of
$\Psi_r=0$ and $\partial\Psi_r/\partial y=0$ gives
\[
 \mu\approx17.5722528999,\qquad y\approx0.004075478609.
\]
Since the endpoint is $a+m-1$, the corresponding limiting
multiplier is approximately $18.5722528999$.
These decimals locate the transition; Proposition~\ref{prop:one-sharp}
proves the uniform bound $18.58$.

The choice $r=2$ can also be justified against every other
integer $r$ within these estimates. For
$r=1,3,4,5,6,7,8$, the exact comparisons
\begin{equation}\label{eq:r-exclusion}
 \left(\frac{163}{60}\right)^{18}
 2^{2L}L^L
 \left(\frac{r^r}{(r+1)^{r+1}}\right)^{18}>1
\end{equation}
give $\Psi_r(18,0)>1$, since $163/60<\ep$. Thus, at $a=18m$,
the logarithmic bound already satisfies
\[
 m\Psi_r(18,0)-1>m-1\ge0.
\]
It cannot give decay, even with the finite-$m$ saving of $1$.
It also fails for smaller $a$, because the coefficient of $a$
in that profile is
$1+\log(r^r/(r+1)^{r+1})<0$.
For $r\ge9$, the degree condition $a>m(2r+1)\ge19m$
itself excludes $a\le18m$.

Thus $r=2$ is the only integer choice that can
reach this scale with the present factorial normalization
and absolute-value decay estimate. The program
\path{checks/optimization.py} reproduces the numerical
transition; \path{checks/endpoints.py} proves
\eqref{eq:r-exclusion} over $\Q$.
